\chapter{Αποτελέσματα και Ανάλυση Δεδομένων}

\section{Αποτελέσματα Ακουστικής Άνεσης στο Μαθησιακό Περιβάλλον}
Η εφαρμογή του Ψηφιακού Διδύμου επέφερε άμεσα και μετρήσιμα αποτελέσματα ως προς την Ποιότητα Εσωτερικού Περιβάλλοντος (\en{IEQ}). Η χρήση του κλειστού βρόχου ελέγχου για τη διατήρηση του θορύβου κάτω από το δυναμικό όριο των 55 dBA αποδείχθηκε ιδιαίτερα αποτελεσματική. Σε αντίθεση με τη χρήση ηχητικών συναγερμών (\en{buzzers}), η οποία θα διατάρασσε την εκπαιδευτική διαδικασία, η διακριτική οπτική ανατροφοδότηση (\en{quiet-lamp}) λειτούργησε ομαλά ως μηχανισμός αυτορρύθμισης της συμπεριφοράς των μαθητών. Τα δεδομένα από την πλατφόρμα Grafana κατέδειξαν σημαντική μείωση στον αριθμό των ημερήσιων υπερβάσεων του ορίου, εξασφαλίζοντας ένα περιβάλλον συμβατό με τις κατευθυντήριες οδηγίες ακουστικής άνεσης.

\begin{figure}[htbp]
\centering
\figureplaceholder{[PLACEHOLDER: ΕΙΚΟΝΑ 6.1 \\ Γράφημα Grafana με τις υπερβάσεις θορύβου \\ ανά ώρα μαθήματος]}
\caption{Απεικόνιση επιπέδων θορύβου και υπερβάσεων (\en{exceedances}) εντός της σχολικής αίθουσας}
\label{fig:noise_results}
\end{figure}

\section{Ενεργειακά Αποτελέσματα και Εξοικονόμηση}
Η εφαρμογή των αυτοματισμών του openHAB, με επίκεντρο τη διαχείριση ενεργοβόρων συσκευών (όπως ο εξοπλισμός του κυλικείου) και την αποτροπή της λειτουργίας «κρυφών» φορτίων (\en{phantom loads}) εκτός ωραρίου, οδήγησε σε εντυπωσιακή εξοικονόμηση. Κατά την περίοδο αναφοράς (Σεπτέμβριος 2024 – Φεβρουάριος 2025), η μέση ημερήσια κατανάλωση της εγκατάστασης ανερχόταν στις 223 kWh.

Κατά τη φάση πλήρους επιχειρησιακής λειτουργίας του Ψηφιακού Διδύμου (Μάρτιος 2025 – Νοέμβριος 2025), η μέση ημερήσια κατανάλωση μειώθηκε στις 179 kWh. Το γεγονός αυτό μεταφράζεται σε μια επαληθευμένη μέση μείωση της τάξης του 19,7\%, η οποία αποτελεί και τον κύριο λειτουργικό δείκτη της εργασίας, καθώς προκύπτει από σταθμισμένη σύγκριση ημερών πλήρους σχολικής λειτουργίας. Η ακατέργαστη σύγκριση σε επίπεδο λογαριασμών, η οποία παρατίθεται αναλυτικά στο Παράρτημα Β, αποτυπώνει ηπιότερη μείωση 10,9\% (201 σε 179 kWh/ημέρα), επιβεβαιώνοντας ότι οι αργίες και οι περίοδοι μειωμένου φορτίου επηρεάζουν ουσιωδώς την ερμηνεία του δείγματος. Αξιοσημείωτη είναι η Year-over-Year (\en{YoY}) σύγκριση για τον μήνα Νοέμβριο: ενώ τον Νοέμβριο του 2024 η κατανάλωση ήταν 252 kWh/ημέρα, τον Νοέμβριο του 2025 κατήλθε στις 170 kWh/ημέρα, σημειώνοντας μείωση 32,5\%, γεγονός που υποδεικνύει την αυξημένη αποδοτικότητα του συστήματος κατά τους χειμερινούς μήνες.
Η τάση αυτή είναι συμβατή με ευρωπαϊκά ευρήματα για ενεργειακή βελτιστοποίηση κτιρίων και δείκτες έξυπνης ετοιμότητας (\en{SRI}) \cite{Apostolopoulos2022, Markoska2019, EPBD2024, Fokaides2020, Zamanidou2024}.

\begin{table}[htbp]
\centering
\caption{Συγκριτική ανάλυση ενεργειακής κατανάλωσης προ και μετά την εφαρμογή του Ψηφιακού Διδύμου}
\label{tab:energy_comparison}
\begin{tabular}{@{}p{5.4cm}ccc@{}}
\toprule
\textbf{Σενάριο Σύγκρισης} & \textbf{Πριν} & \textbf{Μετά} & \textbf{Μεταβολή} \\
\midrule
Ακατέργαστη μέση ημερήσια κατανάλωση λογαριασμών & 201 kWh/ημέρα & 179 kWh/ημέρα & -10,9\% \\
Σταθμισμένη μέση κατανάλωση πλήρους σχολικής λειτουργίας & 223 kWh/ημέρα & 179 kWh/ημέρα & -19,7\% \\
Year-over-Year Νοεμβρίου & 252 kWh/ημέρα & 170 kWh/ημέρα & -32,5\% \\
\bottomrule
\end{tabular}
\end{table}

\section{Αποτελέσματα Τεχνοοικονομικής Ανάλυσης (\en{ROI})}
Βάσει των πραγματικών δεδομένων ενεργειακής εξοικονόμησης και λαμβάνοντας υπόψη μια μέση τιμή ενέργειας της τάξης των 0,162 \euro/kWh, τα οικονομικά μεγέθη της πλατφόρμας είναι εξαιρετικά βιώσιμα. Το συνολικό κόστος της αρχικής επένδυσης (\en{CAPEX}), το οποίο περιλαμβάνει τον εξοπλισμό Edge, τις άδειες και την ανάπτυξη του λογισμικού, υπολογίστηκε στα 19.500 \euro.

Σημειώνεται ότι το αναλυτικό \en{BOM} του Παραρτήματος Α (193,17 \euro) αφορά αποκλειστικά τον πιλοτικό edge εξοπλισμό της εγκατάστασης. Το συνολικό \en{CAPEX} των 19.500 \euro αναφέρεται στο πλήρες έργο επίδειξης και περιλαμβάνει, πέραν του υλικού, την ανάπτυξη λογισμικού, την παραμετροποίηση της υποδομής openHAB/MQTT/InfluxDB, την υλοποίηση της διεπαφής οπτικοποίησης, την εγκατάσταση, το commissioning και την τεχνική υποστήριξη της πιλοτικής λειτουργίας.

\begin{table}[htbp]
\centering
\caption{Συμφιλίωση πιλοτικού \en{BOM} και συνολικού κόστους έργου (\en{CAPEX})}
\label{tab:capex_reconciliation}
\begin{tabular}{@{}p{9.2cm}r@{}}
\toprule
\textbf{Συνιστώσα Κόστους} & \textbf{Ποσό} \\
\midrule
Πιλοτικός edge εξοπλισμός (\en{hardware BOM}, Παράρτημα Α) & 193,17 \euro \\
Ανάπτυξη λογισμικού, παραμετροποίηση πλατφόρμας, εγκατάσταση, commissioning και τεχνική υποστήριξη πιλοτικής λειτουργίας & 19.306,83 \euro \\
\midrule
\textbf{Συνολικό κόστος έργου (\en{CAPEX})} & \textbf{19.500,00 \euro} \\
\bottomrule
\end{tabular}
\end{table}

Ο Πίνακας \ref{tab:capex_reconciliation} καθιστά σαφές ότι το ποσό του Παραρτήματος Α δεν αντιστοιχεί στο πλήρες κόστος του έργου, αλλά μόνο στο υλικό του πιλοτικού κόμβου. Το κύριο τμήμα του \en{CAPEX} αποδίδεται στις υπηρεσίες μηχανικής υλοποίησης, ολοκλήρωσης και επιχειρησιακής θέσης σε λειτουργία της πλατφόρμας.

Στον αντίποδα, τα ετήσια καθαρά οφέλη (\en{OPEX savings}), προερχόμενα από την άμεση εξοικονόμηση ηλεκτρικής ενέργειας και τη βελτιστοποίηση της προγνωστικής συντήρησης, ανέρχονται σε 5.500 \euro. Η ταμειακή αυτή ροή συνεπάγεται περίοδο απόσβεσης (\en{Payback Period}) 42,5 μηνών (περίπου 3,5 έτη). Το συνολικό 5ετές καθαρό κέρδος της επένδυσης εκτιμάται στα 8.000 \euro, παρουσιάζοντας Απόδοση Επένδυσης (\en{5-Year ROI}) της τάξης του 41,02\%.

\begin{figure}[htbp]
\centering
\figureplaceholder{[PLACEHOLDER: ΕΙΚΟΝΑ 6.2 \\ Γράφημα απόσβεσης επένδυσης \\ (\en{Payback Period / ROI})]}
\caption{Διάγραμμα απόσβεσης της επένδυσης και προβολή ταμειακών ροών πενταετίας}
\label{fig:roi_chart}
\end{figure}

